\documentclass[a4paper]{arrowhead}

\usepackage[yyyymmdd]{datetime}
\usepackage{etoolbox}
\usepackage[utf8]{inputenc}
\usepackage[T1]{fontenc}
\usepackage{multirow}
\usepackage{longtable}
\usepackage{tikz}
\usetikzlibrary{positioning,shapes.multipart,arrows.meta}

\renewcommand{\dateseparator}{-}
\setlength{\emergencystretch}{3em}
\lstset{language=HTTP,breakatwhitespace=false,columns=fullflexible}

\newcommand{\fparam}[1]{\textit{\textcolor{ArrowheadBlue}{#1}}}

%% Special references
\newcommand{\mref}[2][]{{\textcolor{ArrowheadPurple}{\hyperref[sec:model:#2]{\ifstrempty{#1}{#2}{#1}}}}}
\newcommand{\oref}[1]{{\textcolor{ArrowheadBlue}{\hyperref[sec:operations:#1]{#1}}}}
\newcommand{\pdef}[1]{{\textcolor{ArrowheadGrey}{#1 \label{sec:model:primitives:#1} \label{sec:model:primitives:#1s}}}}
\newcommand{\pref}[1]{{\textcolor{ArrowheadGrey}{\hyperref[sec:model:primitives:#1]{#1}}}}

% as template_idd.tex; the input and output arguments are passed already formatted (\mref or \pref)
\newrobustcmd\osubsection[5]{
  \addtocounter{subsection}{1}
  \addcontentsline{toc}{subsection}{\protect\numberline{\thesubsection}operation \textcolor{ArrowheadBlue}{#1}}
  \subsection*{
    \thesubsection\quad
    #2 \textcolor{ArrowheadPurple}{#3} \\
    \small
    \hspace*{0.075\textwidth}\begin{minipage}{0.1\textwidth}
      \vspace*{1mm}
      Operation: \\
      \notblank{#4}{Input: \\}{}
      \notblank{#5}{Output: \\}{}
    \end{minipage}
    \begin{minipage}{0.825\textwidth}
      \vspace*{1mm}
      \textcolor{ArrowheadBlue}{#1} \\
      \notblank{#4}{#4 \\}{}
      \notblank{#5}{#5 \\}{}
    \end{minipage}
  }
  \label{sec:operations:#1}
}
\newrobustcmd\msubsection[2]{
  \addtocounter{subsection}{1}
  \addcontentsline{toc}{subsection}{\protect\numberline{\thesubsection}#1 \textcolor{ArrowheadPurple}{#2}}
  \subsection*{\thesubsection\quad#1 \textcolor{ArrowheadPurple}{#2}}
  \label{sec:model:#2} \label{sec:model:#2s}
}
%%

\usepackage{xurl}
\begin{document}
\sloppy

%% Arrowhead Document Properties
\ArrowheadTitle{graphDataHistorian HTTP/TLS/JSON}
\ArrowheadServiceID{graphDataHistorian}
\ArrowheadType{Interface Design Description}
\ArrowheadTypeShort{IDD}
\ArrowheadVersion{5.3.0}
\ArrowheadDate{2026-11-09}
\ArrowheadAuthor{Jerker Delsing}
\ArrowheadStatus{REVIEW DRAFT}
\ArrowheadContact{jerker.delsing@ltu.se}
\ArrowheadFooter{\href{www.arrowhead.eu}{www.arrowhead.eu}}
\ArrowheadSetup
%%

%% Front Page
\begin{center}
  \vspace*{1cm}
  \huge{\arrowtitle}

  \vspace*{0.2cm}
  \LARGE{\arrowtype}
  \vspace*{1cm}

  \Large{Service ID: \textit{"\arrowid"}}

  \vspace*{\fill}

  % Front Page Image

  \vspace*{1cm}
  \vspace*{\fill}

  % Front Page Abstract
  \begin{abstract}
    This document describes the HTTP/1.1 protocol with TLS payload security and JSON payload encoding variant of the \texttt{graphDataHistorian} service of the GraphDBManager system.
  \end{abstract}

  \vspace*{1cm}
\end{center}

\newpage
%%

%% Table of Contents
\tableofcontents
\newpage
%%

\section{Documentation edition and implementation scope}
This is the 5.3.0 documentation edition, dated 2026-11-09. It carries forward the system and service contracts from the retained 5.2.2 documentation and identifies the current primary implementation sources. The former \path|documentation/v5.2.2| folder is retained unchanged.

Document version 5.3.0 is distinct from executable, dependency, registered system/service and wire-contract versions. Values still defined by the source (including 5.2.2 artefact versions, 1.0.0 service versions and \path|SysMonitor-AMQP-IDD-5.2.2| payload identifiers) are preserved. This documentation update does not change runtime behaviour or establish a new implementation release. Historical verification statements describe the tests recorded in the inherited baseline; this revision validates document compilation and references.

\section{Overview}
\label{sec:overview}

This document describes the \texttt{graphDataHistorian} service interface of the GraphDBManager system, which provides persistent storage and query of SenML time series per system and service.

It's implemented using protocol, encoding as stated in below table:
\begin{table}[ht!]
  \centering
  \begin{tabular}{|l|l|l|}
    \hline
    \rowcolor{gray!33} Profile type & Type & Version \\ \hline
    Transfer protocol & HTTP & 1.1 \cite{rfc9112} \\ \hline
    Data encryption & TLS & 1.3 \cite{rfc8446}, 1.2 \cite{rfc5246} \\ \hline
    Encoding & JSON & RFC 8259 \cite{rfc8259} \\ \hline
    Compression & N/A & - \\ \hline
    Semantics & N/A & - \\ \hline
    Ontology & N/A & - \\ \hline
  \end{tabular}
  \caption{Communication and semantics details used for the \texttt{graphDataHistorian} service interface}
  \label{tab:comunication_semantics_profile}
\end{table}

This document provides the Interface Design Description IDD to the \textit{graphDataHistorian -- Service Description} document \cite{gdm-graphDataHistorian-sd}.
For further details about how this service is meant to be used, please consult that document.

The rest of this document describes how to realize the \texttt{graphDataHistorian} service
interface in detail. Both in terms of its operations (Section \ref{sec:operations}) and its information model (Section \ref{sec:model}).

\subsection{Interface registration}

The service is registered in the ServiceRegistry with the interface template \texttt{generic\_\allowbreak{}https}, policy \texttt{CERT\_\allowbreak{}AUTH}, access address \texttt{graphdbmanager}, port 8467 and base path \texttt{/\allowbreak{}graphdbmanager/\allowbreak{}data-\allowbreak{}historian}. The \texttt{operations} property of the interface lists the operations below with their method and path relative to the base path (path parameters are placeholders such as \texttt{\{id\}} that the consumer fills in):

\noindent\begin{longtable}{| p{5.2cm} | p{1.6cm} | p{8.4cm} |} \hline
\rowcolor{gray!33} Operation & Method & Path \\ \hline \endhead
\texttt{list-\allowbreak{}systems} & GET & \texttt{/\allowbreak{}systems} \\ \hline
\texttt{list-\allowbreak{}services} & GET & \texttt{/\allowbreak{}systems/\allowbreak{}\{system\allowbreak{}Name\}/\allowbreak{}services} \\ \hline
\texttt{get-\allowbreak{}data} & GET & \texttt{/\allowbreak{}systems/\allowbreak{}\{system\allowbreak{}Name\}/\allowbreak{}services/\allowbreak{}\{service\allowbreak{}Name\}/\allowbreak{}data} \\ \hline
\texttt{put-\allowbreak{}data} & PUT & \texttt{/\allowbreak{}systems/\allowbreak{}\{system\allowbreak{}Name\}/\allowbreak{}services/\allowbreak{}\{service\allowbreak{}Name\}/\allowbreak{}data} \\ \hline
\end{longtable}

\newpage

\section{Service Operations}
\label{sec:operations}

The interfaces of the \texttt{graphDataHistorian} service, its operations, data models and implementation are provided below.
In particular, each subsection first names the HTTP method and path used to call the operation, after which it names the abstract operation from the \texttt{graphDataHistorian} SD document, as well as input and output types.

All operations require the requester's identity, presented according to the provider's authentication policy: the client certificate (mutual TLS) and, with the \texttt{outsourced} policy, \texttt{Authorization: Bearer IDENTITY-TOKEN//<token>}; systems holding the Authentication secret key send \texttt{Bearer AUTHENTICATOR-KEY//<system>//<hash>}, and with the \texttt{declared} policy \texttt{Bearer SYSTEM//<system>} is sent. Request and response bodies are \texttt{application/json} unless stated otherwise.
Errors are answered with an \mref{ErrorMessage} body and the status codes \texttt{400 Bad Request} (invalid parameter), \texttt{401 Unauthorized} (missing or invalid identity), \texttt{403 Forbidden} (not permitted, e.g. by the management, authorization or blacklist filter), \texttt{404 Not Found}, \texttt{408 Request Timeout}, \texttt{423 Locked}, \texttt{500 Internal Server Error} and \texttt{503 Service Unavailable} (a system it depends on failed).

\osubsection{list-systems}{GET}{/\allowbreak{}graphdbmanager/\allowbreak{}data-historian/\allowbreak{}systems}{}{\mref{SystemListResponse}}

Returns the systems with stored data.

A successful call returns \texttt{200 OK} with the output in the body.

\begin{lstlisting}[label={lst:list-systems},caption={An \oref{list-systems} invocation.}]
GET /graphdbmanager/data-historian/systems HTTP/1.1
Host: graphdbmanager:8467
Authorization: Bearer IDENTITY-TOKEN//<token>
\end{lstlisting}

\begin{lstlisting}[label={lst:list-systems-resp},caption={An \oref{list-systems} response.}]
{
  "systems": [
    "string"
  ]
}
\end{lstlisting}

\osubsection{list-services}{GET}{/\allowbreak{}graphdbmanager/\allowbreak{}data-historian/\allowbreak{}systems/\allowbreak{}\fparam{\{systemName\}}/\allowbreak{}services}{\pref{String}}{\mref{ServiceListResponse}}

Returns the services of a system with stored data.

Path parameter: \fparam{systemName}.

A successful call returns \texttt{200 OK} with the output in the body.

\begin{lstlisting}[label={lst:list-services},caption={An \oref{list-services} invocation.}]
GET /graphdbmanager/data-historian/systems/<systemName>/services HTTP/1.1
Host: graphdbmanager:8467
Authorization: Bearer IDENTITY-TOKEN//<token>
\end{lstlisting}

\begin{lstlisting}[label={lst:list-services-resp},caption={An \oref{list-services} response.}]
{
  "services": [
    "string"
  ]
}
\end{lstlisting}

\osubsection{get-data}{GET}{/\allowbreak{}graphdbmanager/\allowbreak{}data-historian/\allowbreak{}systems/\allowbreak{}\fparam{\{systemName\}}/\allowbreak{}services/\allowbreak{}\fparam{\{serviceName\}}/\allowbreak{}data?\allowbreak{}from=\fparam{\{from\}}\allowbreak{}\&to=\fparam{\{to\}}\allowbreak{}\&count=\fparam{\{count\}}\allowbreak{}\&signals=\fparam{\{signals\}}\allowbreak{}\&signalCount=\fparam{\{signalCount\}}}{\pref{String}}{\pref{Array}$<$\allowbreak{}\mref{SenMLRecord}$>$}

Returns stored SenML records of the system and service, newest first, limited by the query parameters.

Path parameter: \fparam{systemName}, \fparam{serviceName}.

Query parameters: \fparam{from} (Double, optional); \fparam{to} (Double, optional); \fparam{count} (Integer, optional); \fparam{signals} (List\textless{}String\textgreater{}, optional); \fparam{signalCount} (Integer, optional). List values are repeated or comma-separated.

A successful call returns \texttt{200 OK} with the output in the body.

\begin{lstlisting}[label={lst:get-data},caption={An \oref{get-data} invocation.}]
GET /graphdbmanager/data-historian/systems/<systemName>/services/<serviceName>/data?from=<from>&to=<to>&count=<count>&signals=<signals>&signalCount=<signalCount> HTTP/1.1
Host: graphdbmanager:8467
Authorization: Bearer IDENTITY-TOKEN//<token>
\end{lstlisting}

\begin{lstlisting}[label={lst:get-data-resp},caption={An \oref{get-data} response.}]
[
  {
    "bn": "TemperatureSensor/",
    "bt": 1792670400.0,
    "bu": "string",
    "n": "temperature",
    "u": "Cel",
    "v": 21.5,
    "vs": 21.5,
    "vb": true,
    "vd": "string",
    "t": 0.0
  }
]
\end{lstlisting}

\osubsection{put-data}{PUT}{/\allowbreak{}graphdbmanager/\allowbreak{}data-historian/\allowbreak{}systems/\allowbreak{}\fparam{\{systemName\}}/\allowbreak{}services/\allowbreak{}\fparam{\{serviceName\}}/\allowbreak{}data}{\pref{Array}$<$\allowbreak{}\mref{SenMLRecord}$>$}{}

Appends SenML records for the system and service.

Path parameter: \fparam{systemName}, \fparam{serviceName}.

A successful call returns \texttt{204 No Content}.

\begin{lstlisting}[label={lst:put-data},caption={An \oref{put-data} invocation.}]
PUT /graphdbmanager/data-historian/systems/<systemName>/services/<serviceName>/data HTTP/1.1
Host: graphdbmanager:8467
Authorization: Bearer IDENTITY-TOKEN//<token>
Content-Type: application/json

[
  {
    "bn": "TemperatureSensor/",
    "bt": 1792670400.0,
    "bu": "string",
    "n": "temperature",
    "u": "Cel",
    "v": 21.5,
    "vs": 21.5,
    "vb": true,
    "vd": "string",
    "t": 0.0
  }
]
\end{lstlisting}

\newpage

\section{Data Models}
\label{sec:model}

Here, all data objects that can be part of the service calls associated with this service are listed in alphabetic order.
Note that each subsection, which describes one type of object, begins with the \textit{struct} keyword, which is meant to denote a JSON \pref{Object} that must contain certain fields, or names, with values conforming to explicitly named types.
Fields that are not needed may be omitted or \texttt{null} in requests; empty fields are omitted in responses.
As a complement to the primary types defined in this section, there is also a list of secondary types in Section \ref{sec:model:primitives}, which are used to represent things like hashes, identifiers and texts.

\msubsection{struct}{ErrorMessage}




\noindent\begin{tabularx}{\textwidth}{| >{\raggedright\arraybackslash}p{3.9cm} | >{\raggedright\arraybackslash\small}p{5.3cm} | X |} \hline
\rowcolor{gray!33} Object Field & Value Type & Description \\ \hline
"error\allowbreak{}Message" & \pref{String} & Error text, if the event is an error. \\ \hline
"error\allowbreak{}Code" & \pref{Number} &  \\ \hline
"exception\allowbreak{}Type" & \pref{ExceptionType} &  \\ \hline
"origin" & \pref{String} &  \\ \hline
\end{tabularx}



\msubsection{struct}{SenMLRecord}

A SenML record (RFC 8428); a message is a JSON array of records.


\noindent\begin{tabularx}{\textwidth}{| >{\raggedright\arraybackslash}p{3.9cm} | >{\raggedright\arraybackslash\small}p{5.3cm} | X |} \hline
\rowcolor{gray!33} Object Field & Value Type & Description \\ \hline
"bn" & \pref{String} & SenML base name. \\ \hline
"bt" & \pref{Object} & SenML base time (seconds). \\ \hline
"bu" & \pref{String} & SenML base unit. \\ \hline
"bv" & \pref{Object} & SenML base value. \\ \hline
"bs" & \pref{Object} & SenML base sum. \\ \hline
"bver" & \pref{Object} & SenML base version. \\ \hline
"n" & \pref{String} & SenML name. \\ \hline
"u" & \pref{String} & SenML unit. \\ \hline
"v" & \pref{Object} & SenML numeric value. \\ \hline
"vs" & \pref{String} & SenML string value. \\ \hline
"vb" & \pref{Boolean} & SenML Boolean value. \\ \hline
"vd" & \pref{String} & SenML data value (Base64). \\ \hline
"s" & \pref{Object} & SenML sum. \\ \hline
"t" & \pref{Object} & SenML time (seconds, relative to the base time). \\ \hline
"ut" & \pref{Object} & SenML update time (seconds). \\ \hline
\end{tabularx}



\msubsection{struct}{ServiceListResponse}

Service names with data of a system.


\noindent\begin{tabularx}{\textwidth}{| >{\raggedright\arraybackslash}p{3.9cm} | >{\raggedright\arraybackslash\small}p{5.3cm} | X |} \hline
\rowcolor{gray!33} Object Field & Value Type & Description \\ \hline
"services" & \pref{Array}$<$\allowbreak{}\pref{String}$>$ & Service names with stored data. \\ \hline
\end{tabularx}



\msubsection{struct}{SystemListResponse}

List of systems.


\noindent\begin{tabularx}{\textwidth}{| >{\raggedright\arraybackslash}p{3.9cm} | >{\raggedright\arraybackslash\small}p{5.3cm} | X |} \hline
\rowcolor{gray!33} Object Field & Value Type & Description \\ \hline
"systems" & \pref{Array}$<$\allowbreak{}\pref{String}$>$ & System names with stored data. \\ \hline
\end{tabularx}



\subsection{Primitives}
\label{sec:model:primitives}

As all messages are encoded using the JSON format \cite{rfc8259}, the following primitive constructs, part of that standard, become available.
Note that the official standard is defined in terms of parsing rules, while this list only concerns syntactic information.
Furthermore, the \pref{Object} and \pref{Array} types are given optional generic type parameters, which are used in this document to signify when pair values or elements are expected to conform to certain types.

\noindent\begin{tabularx}{\textwidth}{| p{4.6cm} | X |} \hline
\rowcolor{gray!33} Type & Description \\ \hline
\pdef{Array} & A JSON Array \cite{rfc8259}; \pref{Array}$<$A$>$ has elements of type A. \\ \hline
\pdef{Boolean} & JSON Boolean, \texttt{true} or \texttt{false}. \\ \hline
\pdef{ExceptionType} & One of \texttt{ARROWHEAD}, \texttt{INVALID\_\allowbreak{}PARAMETER}, \texttt{AUTH}, \texttt{FORBIDDEN}, \texttt{DATA\_\allowbreak{}NOT\_\allowbreak{}FOUND}, \texttt{TIMEOUT}, \texttt{LOCKED}, \texttt{INTERNAL\_\allowbreak{}SERVER\_\allowbreak{}ERROR}, \texttt{EXTERNAL\_\allowbreak{}SERVER\_\allowbreak{}ERROR} (JSON String). \\ \hline
\pdef{Number} & JSON Number, an integer. \\ \hline
\pdef{Object} & JSON Object; \pref{Object}$<$A$>$ has string keys and values of type A. \\ \hline
\pdef{String} & JSON String. \\ \hline
\end{tabularx}

\newpage

\section{Go and Rust implementation variants}
\subsection{Go implementation}
No Go implementation folder for GraphDBManager is present in the designated primary Go tree. There is no port-specific executable, service binding or parity claim in this edition. A future port needs its own build, state model, security and interface qualification.
\subsection{Rust implementation}
No Rust implementation folder for GraphDBManager is present in the designated primary Rust tree. There is no port-specific executable, service binding or parity claim in this edition. A future port needs its own build, state model, security and interface qualification.

\bibliographystyle{IEEEtran}
\bibliography{bibliography}

\newpage

\section{Revision History}
\subsection{Amendments}

\noindent\begin{tabularx}{\textwidth}{| p{1cm} | p{3cm} | p{2cm} | X | p{4cm} |} \hline
\rowcolor{gray!33} No. & Date & Version & Subject of Amendments & Author \\ \hline

1 & 2026-10-28 & 5.2.2 & Document created from the corresponding DataManager document for the GraphDB-backed copy GraphDBManager. & Jerker Delsing \\ \hline
2 & 2026-10-30 & 5.2.2 & Checked against the primary Java source tree (Eclipse-Arrowhead-java): services, operations, paths and data models unchanged. & Claude Opus 5.5 (draft) \\ \hline
3 & 2026-11-09 & 5.3.0 & New documentation edition; primary-source references and Go/Rust implementation scope; 5.2.2 retained. & Documentation maintenance \\ \hline
\end{tabularx}

\subsection{Quality Assurance}

\noindent\begin{tabularx}{\textwidth}{| p{1cm} | p{3cm} | p{2cm} | X |} \hline
\rowcolor{gray!33} No. & Date & Version & Approved by \\ \hline

1 & & \arrowversion  &  \\ \hline

\end{tabularx}

\end{document}
